\section{问题二：建模与求解}
% TODO: 检测/定位模型（如 YOLO 系）；结构图、损失函数、训练策略、改进点、结果

\subsection{模型结构}
% TODO: 网络结构 TikZ 图（超宽用 \resizebox{\textwidth}{!}{...}）
模型整体结构如图 \ref{fig:arch} 所示。（占位：请补充网络结构图与说明。）

\begin{figure}[H]
\centering
% TODO: 替换为网络结构图，例如 \includegraphics[width=0.9\textwidth]{figures/fig_arch.png}
\caption{模型网络结构}\label{fig:arch}
\end{figure}

\subsection{训练策略与实验设置}
% TODO: 超参数表、实验环境
本文实验环境与超参数如表 \ref{tab:setup} 所示。（占位：请填写配置表。）

\begin{table}[H]
\centering
\zihao{5}
\caption{实验环境与超参数}\label{tab:setup}
\begin{tabular}{llll}
  \toprule
  % TODO: 填写真实配置
  项目 & 配置 & 项目 & 配置 \\
  \midrule
  GPU & -- & 输入尺寸 & -- \\
  框架 & -- & 批大小 & -- \\
  \bottomrule
\end{tabular}
\end{table}

\subsection{改进策略}
% TODO: 数据增强、负样本、损失改进等，未验证项注明「设计/展望」
针对数据特点，本文采用……（占位：请描述改进策略；未验证项务必注明。）

\subsection{训练过程与结果}
% TODO: 训练曲线、PR/混淆矩阵、检测可视化、提交统计
训练过程与验证结果如图表所示。（占位：请补充真实训练曲线与结果。）
